\subsubsection{Learning Temporal Saliency for Time Series Forecasting with Cross-Scale Attention}

\citet{delibasoglu2025crossscale} propose eight synthetic forecasting datasets,
SYN1--SYN8, with predefined relevant feature sets and lag regions. The designs
cover recent contiguous windows, scattered lags, distant windows, and mixtures
of temporal scales. As in the structural-support benchmarks above, the
declared feature--lag locations provide a binary \emph{support} reference:
a location is relevant when used by the synthetic generator. For example, SYN1 assigns
relevance to features 0 and 1 over lags 1--15,
whereas SYN5 assigns relevance to features 1 and 2 over the distant lag window
71--77. The target is constructed from the designated inputs with added
Gaussian noise, which makes the selected regions available as ground-truth
temporal and feature--time saliency. The ground-truth saliency heatmaps in
Figs.~8--9 also show varying color intensity. The binary support inferred from
Table~1 captures membership in the relevant regions, while those plotted colors
contain an additional visual intensity dimension. The paper writes each input
window chronologically and labels its plotted columns $0,\ldots,L-1$, from the
oldest displayed point to the latest. Thus column 0 is an array position at the
old end of the window, whereas the latest input is column $L-1$ and becomes
$\tau=0$ on the common relative-time axis. The paper describes small values in
its separate ``Important Lags'' field as recent; because the synthetic
generation equation and code are absent from the available paper and
repository, an exact algebraic mapping from those lag labels to plot columns
remains unspecified.

\paragraph{Original synthetic-dataset table.}
Table~\ref{tab:crossscale-datasets-exact} reproduces Table~1 of
\citet{delibasoglu2025crossscale}. The source text defines ``Target Noise
level'' as the standard deviation of the Gaussian noise added to the target
$y$. The asterisk attached to SYN8 is retained exactly as displayed in the
source table.

\begin{table}[htbp]
\centering
\small
\setlength{\tabcolsep}{5pt}
\renewcommand{\arraystretch}{1.08}
\caption{Exact transcription of Table~1 in
\citet{delibasoglu2025crossscale}: characteristics of synthetic datasets.}
\label{tab:crossscale-datasets-exact}
\begin{tabular}{@{}
>{\raggedright\arraybackslash}p{0.11\textwidth}
>{\centering\arraybackslash}p{0.45\textwidth}
>{\centering\arraybackslash}p{0.18\textwidth}
>{\centering\arraybackslash}p{0.15\textwidth}@{}}
\toprule
Dataset & Important lags & Important Features & Target Noise level \\
\midrule
SYN1 & 1--15 & $\{0,1\}$ & 0.01 \\
SYN2 & 1--5, 9--10, 15--16, 18, 20, 25--26, 35--36, 50--52, 91--95
     & $\{0,2\}$ & 0.05 \\
SYN3 & 9--10, 15--16, 18, 20--25, 31, 34, 60--65
     & $\{1,2\}$ & 0.08 \\
SYN4 & 9--10, 15--16, 18--21, 41--42, 45--46
     & $\{1,2\}$ & 0.10 \\
SYN5 & 71--77 & $\{1,2\}$ & 0.06 \\
SYN6 & 48--57 & $\{0,2\}$ & 0.05 \\
SYN7 & 60, 62--69 & $\{0,1\}$ & 0.02 \\
SYN8$^{*}$ & mixture of SYN5 and SYN6 & $\{0,1,2\}$ & 0.11 \\
\bottomrule
\end{tabular}
\end{table}

\paragraph{How CrossScaleNet is evaluated.}
On SYN1--SYN8, Figs.~8--9 of \citet{delibasoglu2025crossscale} place the
intrinsic temporal-attention maps of CrossScaleNet and baseline models beside
the known saliency regions. This direct synthetic-ground-truth comparison is
visual, with mask-overlap and ranking scores left untabulated.
Figure~5 additionally shows synthetic feature-importance scores obtained
through post-hoc Feature Ablation and Integrated Gradients. Those plots check
whether the models identify the known important \emph{features}, complementing
the temporal maps, and the synthetic explanation comparison remains graphical.
Forecast quality is measured separately by MSE
and MAE (Fig.~4).

For real Electricity and PM2.5 data, Table~4 reports sufficiency (lower is
better) and comprehensiveness (higher is better) at 10\%, 20\%, and 50\%
selected time-step ratios. These are perturbation-based model-faithfulness
comparisons on datasets without a known generating explanation. Figures~6--7
also compare attention maps visually with post-hoc ShapTime on real examples.
The paper and its official repository provide the reported
sufficiency/comprehensiveness values, while the masking baseline and
forecast-distance calculation remain underspecified for exact
numerical implementation from these sources alone.\footnote{Official
repository linked by the paper: \url{https://github.com/mribrahim/LMS-TSF}.}
Thus, the synthetic reference specifies \emph{where} relevant information
occurs in the input history, while the other tests assess model behavior
without requiring a known generator.

